% answers.tex —— 1977 年·江苏南通地区初试卷 · 参考答案与详细解析片段
% 由 tools/generate_projects.py 依来源总账自 1977/1977_数学真题_江苏南通地区初试_试题.md + 1977/1977_数学真题_江苏南通地区初试_解析正文.md + 1977/1977_数学真题_江苏南通地区初试_答案要点.md 生成，勿手改（改源 MD 后重跑）


\begin{gktip}
固定源 PDF 仅含 1 页题面（原书页码 247），无解答可对照。以下各题均由编者独立推导，最终结论与答案要点逐题一致；原卷大题的中文数字编号已转写为阿拉伯数字，第 1、2 题的小题编号「1．2．」已转写为（1）（2）。
\end{gktip}

\begin{answer}
\gksolhead{1}{}

\textbf{题目}：

回答下列问题：

（1） 当 $a$ 为何值时，$2a>a$；当 $a$ 为何值时，$2a=a$；当 $a$ 为何值时，$2a<a$。

（2） 当 $x$ 为何值时，$|x|=x$。

\Needspace{4\baselineskip}
\textbf{解析}

\gksubpart{（1）}

\Needspace{5\baselineskip}
三问分别解关于 $a$ 的不等式或方程，两边同减 $a$ 即得：

\nopagebreak[4]
\[
2a>a\iff a>0,\qquad
2a=a\iff a=0,\qquad
2a<a\iff a<0.
\]

\gksubpart{（2）}

按绝对值定义：若 $x\ge0$，则 $|x|=x$，等式成立；若 $x<0$，则 $|x|=-x\ne x$，等式不成立。故 $|x|=x$ 当且仅当 $x\ge0$。

\gkanswer{答案（1）：$a>0$ 时 $2a>a$；$a=0$ 时 $2a=a$；$a<0$ 时 $2a<a$。}

\gkanswer{答案（2）：$x\ge0$。}

\end{answer}
\begin{answer}
\gksolhead{2}{}

\textbf{题目}：

写出下列函数的自变量 $x$ 的取值范围（在实数范围内）：

（1） $y=\dfrac{1}{2-x}$；

（2） $y=\sqrt{2x-1}$。

\Needspace{4\baselineskip}
\textbf{解析}

\gksubpart{（1）}

分式有意义要求分母不为零：$2-x\ne0$，即 $x\ne2$。

\gksubpart{（2）}

二次根式有意义要求被开方数非负：$2x-1\ge0$，即 $x\ge\dfrac{1}{2}$。

\gkanswer{答案（1）：$x\ne2$。}

\gkanswer{答案（2）：$x\ge\dfrac{1}{2}$。}

\end{answer}
\begin{answer}
\gksolhead{3}{}

\textbf{题目}：

\Needspace{5\baselineskip}
计算：

\nopagebreak[4]
\[
\left(\frac{2}{3}\right)^{2}+(-5.8)^{0}-\left(\frac{8}{27}\right)^{-\frac{2}{3}}+\left(\frac{1}{2}\right)^{-2}.
\]

\Needspace{4\baselineskip}
\textbf{解析}

\Needspace{5\baselineskip}
逐项计算：

\nopagebreak[4]
\[
\left(\frac{2}{3}\right)^{2}=\frac{4}{9},\qquad
(-5.8)^{0}=1,
\]

\nopagebreak[4]
\[
\left(\frac{8}{27}\right)^{-\frac{2}{3}}
=\left[\left(\frac{8}{27}\right)^{\frac{1}{3}}\right]^{-2}
=\left(\frac{2}{3}\right)^{-2}=\frac{9}{4},\qquad
\left(\frac{1}{2}\right)^{-2}=4.
\]

\Needspace{5\baselineskip}
所以

\nopagebreak[4]
\[
\text{\text{原式}}=\frac{4}{9}+1-\frac{9}{4}+4
=\frac{16+36-81+144}{36}
=\frac{115}{36}.
\]

\gkanswer{答案：$\dfrac{115}{36}$。}

\end{answer}
\begin{answer}
\gksolhead{4}{}

\textbf{题目}：

$k$ 为何值时，方程 $x^{2}-2kx+4=0$ 有两个相等的实数根？并求出这方程的根。

\Needspace{4\baselineskip}
\textbf{解析}

\Needspace{5\baselineskip}
判别式

\nopagebreak[4]
\[
\Delta=(-2k)^{2}-4\times4=4k^{2}-16.
\]

方程有两个相等的实数根当且仅当 $\Delta=0$，即 $k^{2}=4$，得 $k=\pm2$。

当 $k=2$ 时，方程为 $x^{2}-4x+4=0$，即 $(x-2)^{2}=0$，故 $x_{1}=x_{2}=2$；

当 $k=-2$ 时，方程为 $x^{2}+4x+4=0$，即 $(x+2)^{2}=0$，故 $x_{1}=x_{2}=-2$。

\gkanswer{答案：$k=\pm2$；$k=2$ 时 $x_1=x_2=2$，$k=-2$ 时 $x_1=x_2=-2$。}

\end{answer}
\begin{answer}
\gksolhead{5}{}

\textbf{题目}：

已知 $\sin\alpha=\dfrac{3}{5}$（$90^\circ<\alpha<180^\circ$），求 $\cos\alpha$、$\operatorname{tg}\alpha$、$\operatorname{ctg}\alpha$ 的值。

\Needspace{4\baselineskip}
\textbf{解析}

\Needspace{5\baselineskip}
由 $\sin^{2}\alpha+\cos^{2}\alpha=1$，

\nopagebreak[4]
\[
\cos\alpha=\pm\sqrt{1-\left(\frac{3}{5}\right)^{2}}=\pm\frac{4}{5}.
\]

\Needspace{5\baselineskip}
因为 $90^\circ<\alpha<180^\circ$，$\alpha$ 为第二象限角，$\cos\alpha<0$，所以

\nopagebreak[4]
\[
\cos\alpha=-\frac{4}{5}.
\]

\Needspace{5\baselineskip}
从而

\nopagebreak[4]
\[
\operatorname{tg}\alpha=\frac{\sin\alpha}{\cos\alpha}
=\frac{3/5}{-4/5}=-\frac{3}{4},
\qquad
\operatorname{ctg}\alpha=\frac{1}{\operatorname{tg}\alpha}=-\frac{4}{3}.
\]

\gkanswer{答案：$\cos\alpha=-\dfrac{4}{5}$，$\operatorname{tg}\alpha=-\dfrac{3}{4}$，$\operatorname{ctg}\alpha=-\dfrac{4}{3}$。}

\end{answer}
\begin{answer}
\gksolhead{6}{}

\textbf{题目}：

试写出抛物线 $y=1+4x-2x^{2}$ 的顶点坐标、对称轴方程和函数 $y$ 的极值。

\Needspace{4\baselineskip}
\textbf{解析}

\Needspace{5\baselineskip}
配方：

\nopagebreak[4]
\[
y=1+4x-2x^{2}=-2\left(x^{2}-2x\right)+1=-2(x-1)^{2}+3.
\]

抛物线开口向下（二次项系数 $-2<0$），故顶点坐标为 $(1,3)$，对称轴方程为 $x=1$；函数在 $x=1$ 处取得极大值 $3$，无极小值。

\gkanswer{答案：顶点 $(1,3)$，对称轴 $x=1$，$y$ 有极大值 $3$（无极小值）。}

\end{answer}
\begin{answer}
\gksolhead{7}{}

\textbf{题目}：

\Needspace{5\baselineskip}
解方程

\nopagebreak[4]
\[
\frac{\lg 2x}{\lg(x-4)}=2.
\]

\Needspace{4\baselineskip}
\textbf{解析}

\Needspace{5\baselineskip}
定义条件：对数的真数为正且分母不为零，即

\nopagebreak[4]
\[
2x>0,\qquad x-4>0,\qquad \lg(x-4)\ne0,
\]

也就是 $x>4$ 且 $x\ne5$。

\Needspace{5\baselineskip}
去分母并用对数性质：

\nopagebreak[4]
\[
\lg 2x=2\lg(x-4)=\lg\left((x-4)^{2}\right)
\quad\Longrightarrow\quad
2x=(x-4)^{2}.
\]

\Needspace{5\baselineskip}
整理得

\nopagebreak[4]
\[
x^{2}-10x+16=0,
\qquad
(x-2)(x-8)=0,
\]

解得 $x=2$ 或 $x=8$。其中 $x=2$ 不满足 $x>4$，舍去；$x=8$ 满足全部定义条件。故原方程的解为 $x=8$。

\gkanswer{答案：$x=8$。}

\end{answer}
\begin{answer}
\gksolhead{8}{}

\textbf{题目}：

求 200 以内的能被 6 整除的所有自然数的和。

\Needspace{4\baselineskip}
\textbf{解析}

\Needspace{5\baselineskip}
200 以内能被 6 整除的自然数是

\nopagebreak[4]
\[
6\times1,\quad 6\times2,\quad \cdots,\quad 6\times33
\]

\Needspace{5\baselineskip}
（$6\times33=198<200<204=6\times34$）。由等差数列求和公式，它们的和为

\nopagebreak[4]
\[
S=6\times(1+2+\cdots+33)
=6\times\frac{33\times34}{2}
=6\times561=3366.
\]

\gkanswer{答案：$3366$。}

\end{answer}
\begin{answer}
\gksolhead{9}{}

\textbf{题目}：

已知正四棱锥底面边长为 $12\ \mathrm{cm}$，高为 $2\sqrt{7}\ \mathrm{cm}$，求它的侧棱。

\Needspace{4\baselineskip}
\textbf{解析}

正四棱锥的顶点在底面正方形中心的正上方，侧棱连接顶点与底面的顶点。因此高、底面中心到底面顶点的连线（即底面对角线的一半）与侧棱构成直角三角形。

\Needspace{5\baselineskip}
底面正方形边长为 $12$，对角线为 $12\sqrt2$，其一半为 $6\sqrt2$；高 $h=2\sqrt7$。由勾股定理，侧棱长

\nopagebreak[4]
\[
l=\sqrt{h^{2}+\left(6\sqrt2\right)^{2}}
=\sqrt{(2\sqrt7)^{2}+72}
=\sqrt{28+72}=10\ (\mathrm{cm}).
\]

\gkanswer{答案：侧棱长为 $10\ \mathrm{cm}$。}

\end{answer}
\begin{answer}
\gksolhead{10}{}

\textbf{题目}：

两根钢绳悬吊货物时，两条钢绳与铅垂线的夹角分别为 $\alpha$ 和 $\beta$，货物的重量为 $p$，试用表达式分别写出钢绳所承受的拉力 $F_1$ 和 $F_2$。

\Needspace{4\baselineskip}
\textbf{解析}

货物处于平衡状态，所受合力为零。两根钢绳的拉力沿绳方向，与铅垂线的夹角分别为 $\alpha$、$\beta$。

\Needspace{5\baselineskip}
水平方向：

\nopagebreak[4]
\[
F_{1}\sin\alpha=F_{2}\sin\beta;
\]

\Needspace{5\baselineskip}
竖直方向：

\nopagebreak[4]
\[
F_{1}\cos\alpha+F_{2}\cos\beta=p.
\]

\Needspace{5\baselineskip}
由水平方向的方程得 $F_{2}=\dfrac{F_{1}\sin\alpha}{\sin\beta}$，代入竖直方向的方程：

\nopagebreak[4]
\[
F_{1}\left(\cos\alpha+\frac{\sin\alpha\cos\beta}{\sin\beta}\right)=p
\quad\Longrightarrow\quad
F_{1}\cdot\frac{\cos\alpha\sin\beta+\sin\alpha\cos\beta}{\sin\beta}=p,
\]

\Needspace{5\baselineskip}
即

\nopagebreak[4]
\[
F_{1}=\frac{p\sin\beta}{\sin(\alpha+\beta)}.
\]

\Needspace{5\baselineskip}
同理，由对称性（将 $\alpha$ 与 $\beta$ 互换）得

\nopagebreak[4]
\[
F_{2}=\frac{p\sin\alpha}{\sin(\alpha+\beta)}.
\]

\gkanswer{答案：$F_1=\dfrac{p\sin\beta}{\sin(\alpha+\beta)}$，$F_2=\dfrac{p\sin\alpha}{\sin(\alpha+\beta)}$。}

\end{answer}